Table~\ref{tab:planning} gives mission and vehicle success (every weight: supplement S4). The gate routes at most {{VAL:cons_routed_I10_max}} of pairs at 10\,m and {{VAL:cons_routed_I30_max}} at 30\,m: the slip model is uncertain on most dune slopes, so even gentle E1 endpoints are often not TRAVERSABLE on the coarse map ({{VAL:ref_I30_GBM}} of the refusals of boosting at 30\,m and $\tau=0.2$ are at an endpoint). The weight policy routes almost every pair and is chosen for every layer, mostly at the largest weight; with weights up to 1000 (post hoc) held out success at 10\,m rises by up to {{VAL:xw_I10_maxup}} points (E1: interpolation {{VAL:xw_E1_I10_BICUBICcal_up}}, boosting {{VAL:xw_E1_I10_GBM_up}}), and at 30\,m by at most {{VAL:xw_I30_maxup}} for any computable layer.

On 10\,m maps every layer helps: mission success rises from {{VAL:vc_E1_I10_none_m}} with the slip model alone to {{VAL:vc_E1_I10_BICUBICcal_m}} with interpolation ({{VAL:vc_E1_I10_BICUBICcal_md}} points, site level interval {{VAL:vc_E1_I10_BICUBICcal_mci}}), and vehicle success from {{VAL:vc_E1_I10_none_v}} to {{VAL:vc_E1_I10_BICUBICcal_v}}. On 30\,m maps the gain is small: interpolation adds {{VAL:vc_E1_I30_BICUBICcal_md}} mission points (interval {{VAL:vc_E1_I30_BICUBICcal_mci}}) and boosting with Sentinel-2 {{VAL:vc_E1_I30_DEMplusS2_md}} ({{VAL:vc_E1_I30_DEMplusS2_mci}}) to {{VAL:vc_E1_I30_none_m}}, but no computable layer adds more than {{VAL:vc_E1_I30_maxcomp_vd1}} points of vehicle success, and slope, the detector and the network fall below the slip model alone: they flag much of the dune, so routes grow longer and still meet faces. On GLO-30 no layer adds more than {{VAL:vc_E1_G30_maxcomp_md1}} points, on AW3D30 interpolation adds {{VAL:prp_A30_d1}} and on FABDEM {{VAL:prp_F30_d1}} (supplement Table~S13). The sensitivity vehicle and pair distance change little (Tables~S10 and~S14). With E2 endpoints the slip model alone reaches {{VAL:vc_E2_I30_none_m}} at 30\,m and the gain of interpolation is {{VAL:vc_E2_I30_BICUBICcal_md}} (interval {{VAL:vc_E2_I30_BICUBICcal_mci}}).

The face map oracle reaches {{VAL:vc_E1_I30_FACEMAP_m}} mission success on the 30\,m grid but only {{VAL:vc_E1_I30_FACEMAP_v}} vehicle success: {{VAL:vc_E1_I30_FACEMAP_feas}} of its successes are drivable, because its routes go round faces over slopes the coarse map shows as gentle. Registered degraded face maps at $w=5$ (supplement Table~S7) show how exact face information must be: the oracle gives {{VAL:dg_E1_I30_FACEMAP_risk}} at this weight, displaced by 10\,m {{VAL:dg_E1_I30_shift10_risk}} and by 30\,m {{VAL:dg_E1_I30_shift30_risk}}, and maps from simulated 2\,m and 5\,m surveys only {{VAL:dg_E1_I30_survey2_risk}} and {{VAL:dg_E1_I30_survey5_risk}}, because the 28\,degree rule misses faces on a coarser grid.
